The equality \(Q^{CR}_{\bar P_0,C_C}=Q^{CR}_{P_0,C_C}\) follows assignment by assignment because the two schedule laws have the same marginal \(P_{0,z}\) for every observed \(z\).  Conditional on the fixed-count label vector, the cluster observations are independent.  The same differentiation as in (12) gives derivative \(h_k\) for \(U_k\), and the conditional observed score is \(\dot\ell_{c,k}^{CR}\).  It is centered within stratum \(k\), and
\[
 \frac1C\sum_{c=1}^C\mathbb E[\dot\ell_{c,j}^{CR}\dot\ell_{c,k}^{CR}\mid L]
 ={\mathbf 1}\{j=k\}\frac{C_{C,k}}C\frac1{M_k\tau_k^4}
 \sum_{z\in\mathcal S_k}\operatorname{Var}_{P_0}(W(z))
 \longrightarrow{\mathbf 1}\{j=k\}\frac{\alpha_k}{\tau_k^2}. \tag{16}
\]
This is \(I_{CR}\).  Applying the bounded log-normalizer expansion (14) to the fixed strata makes the linear term \(h'\Delta_C^{CR}\), the quadratic term converge to \(-h'I_{CR}h/2\), and bounds the summed cubic remainder by \(O(C^{-1/2})\) on bounded \(h\).  The maximum normalized score is \(O(C^{-1/2})\), while (16) gives its total covariance.  Expanding the product of conditional characteristic functions therefore proves \(\Delta_C^{CR}\Rightarrow N(0,I_{CR})\) and the asserted LAN expansion.

The triangular-array estimator expansion (8), applied under \(\bar P_{C,h}\), gives
\[
 \sqrt C\{\widetilde U_{C,A}-U_A(P_0)\}\Rightarrow
 (h_k+\sqrt{\tau_k^2/\alpha_k}Z_k)_{k\in A}. \tag{17}
\]
For any finite prior on bounded points of \(H_A\), joint convergence and uniform integrability of the likelihood ratios imply convergence of its Bayes decision problem to the Gaussian problem.  The compact-ball minimax and finite-prior approximation argument given in the Bernoulli proof applies verbatim because the action set is finite and all active variances in (16) are positive.  Letting the ball radius increase proves the displayed lower bound with \(G_A(\tau_A^2,\alpha_A)\).  Label deletion does not change this benchmark experiment.

For the final witness, let each target slice reveal an independent bit with law \(\operatorname{Bernoulli}(1/2)\), constant across the \(n\) unit outcomes for a realized target assignment.  Then \(W\) equals that bit, so \(\tau_k^2=1/4\).  Each of the asymptotic fraction \(\alpha_k\) of observations in stratum \(k\) has Bernoulli mean information \((1/4)^{-1}=4\); hence the per-cluster information is \(4\alpha_k\), as claimed.